\section{An explicit choice of the constant in K17-VK-EXP7}
\label{sec:kt-c7-explicit}

\paragraph{Finding KT-C7-EXPLICIT.}
Under the stipulated archive theorems and the carried K17-VK-PT and
K17-VK-EXP7 argument, one explicit admissible choice is
\[
 C_7=10^{\,10^{3500}},\qquad
 \log|K_{4k}(\mathbf Z)|\le C_7(2k)^7\quad(k\ge1).
\]
This is a very coarse upper bound for a permissible constant.  It is
not the least constant and offers no numerical vanishing criterion.
Its purpose is to remove the unevaluated cohomology constants from the
source reconstruction requested in item 4.

All source locators below are relative to
\path{src/Development/AutomorphicArithmetic/AutomorphicVarieties/}
at commit \path{7cbf7a0c32fca1badf05931879a37c3f8a2b03b4}.
The theorem \path{cohBoundFree},
\path{AutomorphicForms/TemperedLocalComponents/CohBoundFree.lean}:93--99,
returns existential constants.  Its helper \path{cohBound_of}:69--91
obtains a base-level witness $c$ at line 75 and returns
$(c3^{dn^2},dn^2)$ at line 79, with $d=[F:\mathbf Q]$.
The witness ultimately comes from finite orbits and stabilizer
cohomology.  The following retention of the source's explicit lower
bounds replaces it with an integer expression.

\subsection{An explicit arithmetic-group cohomology bound}

Let $F$ be imaginary quadratic of absolute discriminant $D$ and let
$n\ge1$.  Define
\[
 I=D^n(2n)^{2n},\quad R=n^nI^2,\quad T=(2R+1)^{2n},
\]
\[
 A=I^{1+2n^2},\quad B=(T+1)^{2n^2},\quad H=AB^2T^n,
 \qquad b(D,n)=HAB^H.
\tag{KT-C7.1}
\]
\paragraph{Finding KT-COH-EXPLICIT.}
For every finite-dimensional complex representation $V$ of
$\mathrm{GL}_n(\mathcal O_F)$,
\[
 \sum_{q\ge0}\dim H^q(\mathrm{GL}_n(\mathcal O_F),V)
 \le b(D,n)\dim V.
\tag{KT-C7.2}
\]

\paragraph{Proof.}
\path{HermiteNormalForms/SmulStrictMono.lean}:17--19,244--249 defines
$\mathrm{dd}=dn$, $I_0=D^n(dn)^{dn}$, $RR=n^nI_0^2$, and
\[
 \operatorname{orbBound}(q)
 =I_0^{1+dn^2}\bigl((2RR+1)^{dn}+1\bigr)^{dn^2(q+1)}.
\]
For $d=2$ this is exactly $AB^{q+1}$.
\path{HermiteNormalForms/ColsMinSmul.lean}:34--75 bounds the number
of group orbits of ordered $q$-simplices of its Hermitian nerve by
this integer.  The augmented permutation resolution is exact by the
archive theorem \path{ConeExtensionAcyclic.sx_exact}, used at
\path{BuildingsResidueChars/BoundedComplexOfBoundedSimplicialComplex.lean}:18.

First bound a vertex stabilizer.  By
\path{HermiteNormalForms/Stabfin.lean}:24--26 choose a nonsingular
matrix $A_0$ whose $n$ columns are minimal vectors of the vertex $v$.
Every minimal vector of $v$ is in the set
$\operatorname{nearVec}(A_0)$ defined at
\path{HermiteNormalForms/SmulStrictMono.lean}:258--259.
By \path{HermiteNormalForms/Nearvec.lean}:9--38 this set has at most
$T$ elements.  A stabilizer element is determined by the images of
the $n$ columns: equality on those columns gives equality after
multiplication by $A_0$, and then equality because $\det A_0\ne0$.
This is the injection in \path{Stabfin.lean}:35--50, now with its
cardinality retained.  Hence
\[
 |\operatorname{Stab}(v)|\le T^n.
\tag{KT-C7.3}
\]

Choose one representative from each orbit of ordered $1$-simplices;
there are at most $AB^2$.  The injection in
\path{BuildingsResidueChars/AdelicRepDimensionBounded.lean}:22--68
maps the $q+1$ distinct vertices of an injective ordered $q$-simplex
into the disjoint union of the stabilizers of the first vertices of
these representatives.  This union has at most $AB^2T^n=H$ elements.
There is consequently no injective $q$-simplex when $q\ge H$.

The original resolution includes repeated tuples.  The normalized
alternating endomorphism is a chain map by
\path{BuildingsResidueChars/AltChainOfPermSplit.lean}:5--46.
It is the identity on the augmentation and zero in a degree in which
all tuples repeat a vertex, by
\path{BuildingsResidueChars/RepHomAlternatingSumEqZero.lean}:10--57.
The projective-resolution comparison in
\path{BuildingsResidueChars/ExtVanishingOfResolutionEndo.lean}:32--38
therefore gives $H^q=0$ for $q\ge H$.
Projectivity follows from finite stabilizers over $\mathbf C$, exactly
as in \path{AutomorphicForms/TemperedLocalComponents/BddArith.lean}:41--49.

In degree $q$, an equivariant homomorphism from the permutation term
to $V$ is determined by one vector for each orbit.  Its dimension is
at most $AB^{q+1}\dim V$; this is the source argument
\path{BuildingsResidueChars/PermutationModuleHomologyRank.lean}:4--34.
Cohomology is a subquotient of this Hom term.  Consequently
\[
 \sum_q\dim H^q
 \le\left(\sum_{q=0}^{H-1}AB^{q+1}\right)\dim V
 \le HAB^H\dim V.
\]
This proves KT-COH-EXPLICIT. \hfill$\square$

\subsection{The principal-level coefficient}

If $F$ has class number one, write $K(m)$ for the principal subgroup
of $\mathrm{GL}_n(\widehat{\mathcal O}_F)$ of level $m$.  We claim that
the same total adelic cohomology quantity used by \path{cohBoundFree}
satisfies
\[
 \operatorname{cohDim}(\mathrm{GL}_n(F),K(m),\xi)
 \le b(D,n)m^{2n^2}\dim\xi.
\tag{KT-C7.4}
\]

An adelic matrix $g$ determines a full $\mathcal O_F$-lattice $L$ in
$F^n$ whose completion is $g\widehat{\mathcal O}_F^n$.
For an explicit construction, choose an integer $a>0$ with
\[
 a\widehat{\mathcal O}_F^n\subseteq
 g\widehat{\mathcal O}_F^n\subseteq
 a^{-1}\widehat{\mathcal O}_F^n.
\]
The middle lattice gives a submodule of the finite quotient
$a^{-1}\mathcal O_F^n/a\mathcal O_F^n$ by Chinese remainders; take
its inverse image for $L$.  The same finite quotient shows that its
completion is the prescribed one.  Since $\mathcal O_F$ is a PID,
$L=\gamma\mathcal O_F^n$ for a matrix $\gamma\in\mathrm{GL}_n(F)$.
Thus $g=\gamma k$ with $k\in K(1)$.
There is one rational orbit at level one, with stabilizer
$\mathrm{GL}_n(\mathcal O_F)$.

The fibers at level $m$ have cardinality
$[K(1):K(m)]\le m^{2n^2}$, by
\path{BuildingsResidueChars/LevelSubgroupIndexBoundOfCardOhat.lean}:22--38.
The Shapiro fiber argument
\path{BuildingsResidueChars/OrbitBoundOfShapiroFunctorial.lean}:3--23,
with KT-COH-EXPLICIT, proves (KT-C7.4).
This is the one-orbit, base-level-one specialization of the helper
\path{gen} in \path{CohBoundFree.lean}:28--60.

\subsection{Substituting six explicit fields}

Take $F_0=\mathbf Q(\sqrt{-3})$ for the rank-three branch and take
the six rank-two fields with absolute discriminants
\[
 \mathcal D=\{3,7,11,19,43,67\}.
\]
They are pairwise coprime, so they satisfy the same six-field
requirement as K17's illustrative choices.
All have class number one.  To check this, Minkowski's bound says
that every ideal class has an integral representative of norm at
most $(2/\pi)\sqrt D$ (J.~S.~Milne,
\emph{Algebraic Number Theory}, version 3.08, Theorem 4.3, p.~70;
\url{https://www.jmilne.org/math/CourseNotes/ANT.pdf}).
For $D=3,7$ this is $<2$.  For $D=11,19$ it is $<3$, and $2$ is
inert.  For $D=43$ it is $<5$, and $2,3$ are inert.  For $D=67$ it
is $<6$, and $2,3,5$ are inert.  In the last two cases the only
possible nonunit norm is $4$, whose ideal is the principal ideal
$(2)$.  Thus all representatives are principal.  Inertness follows
from $-D\equiv5\pmod8$ at $2$ and the indicated quadratic
nonresidues at the odd primes; the accompanying exact script checks
each instance.

K17, \path{evidence/work/k_quantitative_refinement.md},
sections 4--5, permits constants $B_0$ for rank three at level $6$
and $B_D$ for rank two at level $18$.  Hence (KT-C7.4) allows
\[
 B_0=6^{18}b(3,3),\qquad B_D=18^8b(D,2).
\tag{KT-C7.5}
\]
The levels and the use of the full rational group are explicit in
\path{AutomorphicForms/TemperedLocalComponents/KlevThree.lean}:18--28
and \path{GaloisLibrary/CrystallineBaseChange/KlevThree.lean}:17--27:
both use \path{cohBoundFree ... n top} and replace a level $M$ by
$3M$.  The carried K17 exponent-seven proof therefore permits the
fully explicit integer
\[
 C_7^*=20\left(6^{18}b(3,3)+32\,18^8
                 \sum_{D\in\mathcal D}b(D,2)\right).
\tag{KT-C7.6}
\]
Every quantity in this expression is now given by integer additions,
multiplications, and powers in (KT-C7.1).

\subsection{Evaluation by elementary decimal bounds}

For $(D,n)=(3,3)$, $I=1{,}259{,}712<10^7$ and $2R+1<10^{14}$.
Thus
\[
 T<10^{84},\quad A<10^{133},\quad B\le10^{1512},
 \quad H<10^{3409},\quad \log_{10}b(3,3)<10^{3413}.
\]
For $n=2$ and $D\le67$, $I\le1{,}149{,}184<10^7$ and again
$2R+1<10^{14}$.  Hence
\[
 T<10^{56},\quad A<10^{63},\quad B\le10^{448},
 \quad H<10^{1071},\quad \log_{10}b(D,2)<10^{1074}.
\]
The $+1$ in $B$ causes no loss: an integer $T<10^u$ has
$T+1\le10^u$.
Also $6^{18}<10^{15}$, $18^8<10^{11}$, and
$20(10^{15}+192\cdot10^{11})<10^{18}$.
It follows that
\[
 C_7^*<10^{18+10^{3413}}<10^{10^{3414}}<10^{10^{3500}}.
\]
This proves KT-C7-EXPLICIT.  The exact-arithmetic replay is
\path{check_c7_arithmetic.py}; it verifies these inequalities without
trying to expand a double-exponential integer.

\paragraph{Scope.}
The source-standing result and the constant calculation are distinct.
The former is an explicit substitution in named archive theorems;
the latter is a new elementary quantitative refinement of the source's
finite-orbit proof, using those theorems as inputs.  Neither is a new
Lean kernel certificate.  The geometric constants arise from
Hermitian reduction complexes and finite congruence-level fibers.
Their interpretation as a bound for integral Tate pieces over
$\operatorname{Spec}\mathbf Z$ uses the carried K17-VK-PT bridge.
They do not establish vanishing or Vandiver's conjecture.
